\section{Research Methodology}
\label{sec:Res}

This section describes the empirical setup used to evaluate a minimal runtime verification layer over an automotive ECU telemetry stream. The approach is \textbf{trace-driven} and \textbf{monitor-centric}: recorded ECU time-series are replayed server-side, and an external runtime verifier checks those streams against formalised properties. The design follows a quasi-experimental logic~\cite{saunders2016research}: replay traces act as controlled behaviour sources, while monitors provide repeatable verdicts on normality versus violation. The deployment assumption is that an IoT/ECU endpoint forwards telemetry records to a server-side analysis service; hence measured overhead reflects server-side monitoring cost for ingested records, not direct on-device ECU execution overhead. The project artefacts, scripts, and implementation used in this study are available in the public repository: \url{https://github.com/Naronitsu/RD2}.

\subsection{Overall approach}

The methodology proceeds in four stages: (1)~prepare a cleaned ECU dataset with stable schema and timestamps (sourced from the Kaggle Car ECU Datalogs dataset~\cite{gxkok21ecudatalogs}); (2)~construct an abnormal variant by injecting controlled violations into existing rows while preserving file size; (3)~deploy \textbf{LARVA}~\cite{colombo2009larva} in a Spring-based server endpoint that processes each sample and raises monitor states; and (4)~benchmark latency and throughput by replaying the same input under RV enabled versus disabled execution paths.

Prior integration work shows LARVA can be composed with event middleware pipelines~\cite{abela2023rvtee}. In this study, a server-side CSV replay service supplies the ordered event stream (analogous to middleware or bus traffic), and LARVA acts as the oracle without requiring ROS.

\subsection{Instrumentation: LARVA and flags}

LARVA supports specification-driven monitoring of timed behaviour~\cite{colombo2009larva}. Each \emph{flag} below is implemented as a monitor property that raises a violation when its predicate becomes false. The evaluated assertion flags are:

\begin{enumerate}
    \item \textbf{Voltage out of range under running RPM.} Trigger when RPM is above idle threshold and battery voltage is outside calibrated bounds.
    \item \textbf{Coolant too high.} Trigger when coolant temperature exceeds a configured thermal safety limit.
    \item \textbf{Throttle spike.} Trigger when throttle delta exceeds a threshold within a short runtime interval.
    \item \textbf{Load--MAP mismatch.} Trigger when absolute disagreement between Load and MAPSource exceeds tolerance.
\end{enumerate}

Each flag is evaluated on replayed traces; the abnormal dataset injects approximately 20\% violation rows distributed across the four flags to validate that the corresponding LARVA specification fires, while benign traces provide a baseline for false-positive checks.

\subsection{Benchmark protocol}

Two benchmark regimes are executed with identical machine and JVM conditions: \textbf{baseline} (clean trace) and \textbf{abnormal} (injected trace). In both regimes, the same endpoint and CSV shape are used, and only RV instrumentation is toggled through a per-request parameter (\texttt{rvEnabled=true/false}) and the global server setting (\texttt{rv.enabled}). The primary overhead summary uses the stock generated monitor with one warmup iteration and five measured trials per mode; supplementary sessions (e.g., threshold-tuning repeatability) use three warmup iterations and ten measured trials per mode. Reported metrics include rows processed, elapsed time, throughput (rows/s), average, standard deviation, p95/p99 latency, and overhead. Each elapsed-time value represents end-to-end processing of the full input trace per request (68{,}433 rows), not per-row latency:
\[
\text{overhead}(\%)=\frac{\overline{t}_{\mathrm{RV\ on}}-\overline{t}_{\mathrm{RV\ off}}}{\overline{t}_{\mathrm{RV\ off}}}\times 100.
\]

\subsection{Research questions and validation}

The following questions guide evaluation: (RQ1)~Do the four selected properties detect the injected abnormal classes with high recall? (RQ2)~What is the observed monitoring overhead relative to non-instrumented processing on the same server path? (RQ3)~How stable are latency and throughput across repeated runs (mean, standard deviation, and tail percentiles)? Answers are sought by structured log analysis and repeated trials, consistent with cybersecurity V\&V emphasis on repeatable test cases~\cite{ekert2021cybersecurity}.

\subsection{Limitations}

Results apply to the simulation and property set instantiated here; they do not replace homologation testing or formal verification of production ECU binaries. Future work could couple the same monitors with recorded CAN traces or hardware-in-the-loop benches.

\subsection{Ethical considerations}

This study uses publicly available telemetry data and synthetic fault injection; it does not involve human participants, personal identifiers, or intervention on live vehicle systems. Dataset handling is limited to technical signal analysis for research purposes, and all transformations are reproducible through version-controlled scripts. To support research integrity and transparency, benchmark procedures, analysis code, and artefacts are disclosed in the project repository, enabling independent replication and scrutiny.
